\documentclass{article}
\usepackage{mathptmx}
\usepackage{helvet}
\usepackage{courier}
\usepackage[T1]{fontenc}
\usepackage[latin9]{inputenc}
\usepackage{array}
\usepackage{url}
\usepackage{amsmath}
%
\usepackage{setspace}
%\onehalfspacing
\usepackage{amssymb}
\usepackage{graphicx}                              %for PNG images (pdflatex)
%\usepackage{graphics}                              %for EPS images (latex)
\usepackage[linkbordercolor={1.0 1.0 0.0}]{hyperref} %for \url tag
\usepackage{color}                                 %for defining custom colors
\usepackage{framed}                                %for shaded and framed paragraphs
\usepackage{textcomp}                              %for various symbols, e.g. Registered Mark
\usepackage{geometry}                              %for defining page size
\usepackage{longtable}                             %for breaking tables
%
\geometry{verbose,a4paper,tmargin=2.5cm,bmargin=2.5cm,lmargin=2.5cm,rmargin=2cm}
\hypersetup{
  pdfauthor = {D.Cameron, A.Konstantinov},
  pdftitle = {The NorduGrid GridFTP Server.},
  pdfsubject = {},
  pdfkeywords = {},
  pdfcreator = {PDFLaTeX with hyperref package},
  pdfproducer = {PDFLaTeX}
}
%
\bibliographystyle{IEEEtran}                       %a nice bibliography style
%
\def\efill{\hfill\nopagebreak}%
\hyphenation{Nordu-Grid}
\setlength{\parindent}{0cm}
\setlength{\FrameRule}{1pt}
\setlength{\FrameSep}{8pt}
\addtolength{\parskip}{5pt}
\renewcommand{\thefootnote}{\fnsymbol{footnote}}
\renewcommand{\arraystretch}{1.3}
\newcommand{\dothis}{\colorbox{shadecolor}}
\newcommand{\globus}{Globus Toolkit\textsuperscript{\textregistered}~2~}
\newcommand{\GT}{Globus Toolkit\textsuperscript{\textregistered}}
\newcommand{\ngdl}{\url{http://ftp.nordugrid.org/download}~}
\definecolor{shadecolor}{rgb}{1,1,0.6}
\definecolor{salmon}{rgb}{1,0.9,1}
\definecolor{bordeaux}{rgb}{0.75,0.,0.}
\definecolor{cyan}{rgb}{0,1,1}
%
\begin{document}
\def\today{\number\day/\number\month/\number\year}

\begin{titlepage}

\begin{tabular}{rl}

\resizebox*{3cm}{!}{\includegraphics{ng-logo.png}}

&\parbox[b]{2cm}{\textbf \it {\hspace*{-1.5cm}NORDUGRID\vspace*{0.5cm}}}

\end{tabular}

\hrulefill

{\raggedleft NORDUGRID-TECH-26\par}

{\raggedleft \today\par}

\vspace*{2cm}

%%%%---- The title ----
{\centering \textsc{\Large The NorduGrid GridFTP Server}\Large \par}
\vspace*{0.5cm}

%%%%---- A subtitle, if necessary ----
{\centering \textit{\large Description and Administrator's Manual}\large \par}

\vspace*{1.5cm}
%%%%---- A list of authors ----
{\centering \large A. Konstantinov\footnote{aleks@fys.uio.no}, D. Cameron \large \par}

\end{titlepage}

\tableofcontents

\newpage

\section{Introduction\label{sec:intro}}

The NorduGrid~\cite{nordugrid} GridFTP service (GFS) has historically
been the gateway between Grid users and their jobs running on
ARC-enabled resources. It provides an interface for secure job
submission and retrieval with authorization and authentication based
on the Grid Security Infrastructure~\cite{gsi}, and consists of a
standard GridFTP server with NorduGrid modifications on top to handle
Grid jobs. These jobs are then processed and sent to the local batch
system by the Grid Manager (GM).

With the evolution of Grid technologies, and the need for more
standardized interfaces to allow different Grid products to
communicate with each other, the Hosting Environment Daemon
(HED)~\cite{HED} framework and associated Advanced Resource Execution
Service (A-REX) web service were developed to replace the aging GFS
and GM with a modern standards-compliant product. However, the
transition to a new framework takes time and so the legacy GFS is
still supported and maintained until such time as it becomes
obsolete. This document is designed for administrators who wish to run
the GFS alongside A-REX and describes how to configure and run the GFS
service. For installation of the GFS, up to date instructions may be
found on the NorduGrid website \textbf{\url{http://www.nordugrid.org}}.

\section{Main Concepts}

A detailed summary of job workflow in ARC may be found in the A-REX
technical documentation~\cite{arex}. The following instructions
assume knowledge of concepts such as the \emph{session directory},
job states and how data management is performed by A-REX.

The GFS provides a means to map GSI identities to local usernames, and
thus can expose a local filesystem to the Grid using a highly
configurable set of authentication policies. It also allows Grid job
submission by providing a way for clients to upload a job description
to a Grid resource, and for that job to be executed on a resource
under a mapped local username.

Local file access in the GFS is implemented through plugins (shared
libraries). There are 3 plugins provided: \textit{fileplugin.so,
  gaclplugin.so} and \textit{jobplugin.so} . The
\textit{fileplugin.so} is intended to be used for plain file access
with the configuration sensitive to the user subject and is not
necessary for setting up a NorduGrid compatible site. The
\textit{gaclplugin.so} uses GACL~\cite{gacl} to control access to the
local file system. The \textit{jobplugin.so} uses information about
jobs being controlled by A-REX and provides access to session
directories of the jobs owned by the user. It also provides an
interface (virtual directory and virtual operations) to submit,
cancel, clean, renew credentials and obtain information about the job.

To make GFS to inter-operate with other parts of ARC only one
\emph{jobplugin.so} needs to be configured.


\section{Configuration}

The GFS configuration is done through a single INI-style configuration
file. The XML-style configuration supported by A-REX are not supported
by the GFS, and so care must be taken if using the GFS with A-REX
configured using XML-style configuration. The safest option is to use
the same INI-style file for A-REX and the GFS. The default location of
the GFS configuration file is

\begin{itemize}
\item \textit{/etc/arc.conf}
\end{itemize}

A different configuration file location can be specified by the
environment variable ARC\_CONFIG. The configuration file
consists of empty lines, lines containing comments (lines starting
with \#) or configuration commands. It is separated into
sections. Each section starts with a string containing

\begin{itemize}
\item \emph{{[}section name/subsection name/subsubsection name]}. 
\end{itemize}

Each section continues until the next section or until the end of the
file. The configuration file can have commands for multiple
services/modules/programs. Each service has its own section named
after it. The GFS uses the \emph{{[}gridftpd]} section and
sub-sections, along with other authorization-related
sections. Commands in section \emph{{[}common]} apply to all services
configured in the configuration file. Command lines have the format

\begin{itemize}
\item \emph{name=''arguments string''.}
\end{itemize}

\subsection{General Configuration Parameters}

The following parameters are defined in the \emph{{[}gridftpd]}
section of the configuration file.

\begin{itemize}
\item \textbf{\textit{pidfile}}\textit{=path} -- specifies file where
  process id of the gridftpd process will be stored. Defaults to
  \emph{/var/run/gridftpd.pid} if running as root and
  \emph{\$HOME/gridftpd.pid} otherwise.
\item \textbf{\textit{logfile}}\textit{=path} -- specifies name of file
  for logging debug/informational output. Defaults to
  \emph{/var/log/arc/gridftpd.log}. If installed from packages, the
  default log is managed by logrotate.
\item \textbf{\textit{logsize}}\textit{=size number} -- restricts
  log file size to \emph{size} and keeps \emph{number} archived log
  files.
\item \textbf{\textit{debug}}\textit{=number} -- specifies level of
  debug information. More information is printed for higher
  levels. Currently the highest effective number is 5 (DEBUG) and
  lowest 0 (FATAL). Defaults to 3 (INFO).
\item \textbf{\textit{port}}\textit{=number} - specifies TCP/IP port
  number.  Default is 2811.
\item \textbf{\textit{include}}\textit{=path} - include contents of
  another file. Generic commands cannot be specified there.
\item \textbf{\textit{encryption}}\textit{=yes|no} - specifies if
  server will allow data transfer to be encrypted. Default is yes.
\item \textbf{\textit{pluginpath}}\textit{=path} - specifies the path
  where plugin libraries are installed. In a normal installation this
  is \emph{\$(ARC\_LOCATION)/lib/arc}.
\item \textbf{\textit{allowunknown}}\textit{=yes|no} - if set to
  \emph{yes}, clients are not checked against the grid-mapfile. Hence
  only access rules specified in this configuration file will be
  applied.
\item \textbf{\textit{firewall}}\textit{=hostname} - use IP address of
  the \textit{hostname} in response to PASV command instead of IP
  address of a network interface of the computer. An IP address can be
  used instead of \textit{hostname}. This command may be useful if the
  server is situated behind a NAT.
\item \textbf{\textit{unixgroup}}\textit{=group rule} - define local
  UNIX user and optionally UNIX group to which user belonging to
  specified authorization \textit{group} is mapped (see
  Section~\ref{sec:Authorization} for definition of group). Local
  names are obtained from the specified \textit{rule}. If the
  specified rule could not produce any mapping, the next command is
  used. Mapping stops at first matched rule. The following rules are
  supported:

\begin{itemize}
\item \textbf{\textit{mapfile}} \textit{file} - the user's subject is
  matched against a list of subjects stored in the specified file, one
  per line followed by a local UNIX name.
\item \textbf{\textit{simplepool}} \textit{directory} - the user is
  assigned one of the local UNIX names stored in a file
  \textit{directory/pool}, one per line. Used names are stored in
  other files placed in the same \textit{directory}. If a UNIX name
  was not used for 10 days, it may be reassigned to another user.
\item \textbf{\textit{lcmaps}} \textit{library directory database} -
  call LCMAPS functions to do mapping. Here \textit{library} is the
  path to the shared library of LCMAPS, either absolute or relative to
  \textit{directory}; \textit{directory} is the path to the LCMAPS
  installation directory, equivalent to the LCMAPS\_DIR variable;
  \textit{database} is the path to the LCMAPS database, equivalent to
  the LCMAPS\_DB\_FILE variable. Each argument except \textit{library}
  is optional and may be either skipped or replaced with '{*}'. It is
  important to ensure that no configured LCMAPS plugin performs switch
  of local user identity (setuid). That may interfere with way the GFS
  handles local user identities.
\item \textbf{\textit{mapplugin}} \textit{timeout plugin}
  {[}\textit{arg1} {[}\textit{arg2} {[}...]]] - run external
  \emph{plugin} executable with specified arguments. Execution of
  \emph{plugin} may not last longer than \emph{timeout} seconds. A
  rule matches if the exit code is 0 and there is a UNIX name printed
  on \emph{stdout}. A name may be optionally followed by a UNIX group
  separated by ':'. In arguments the following substitutions are
  applied before the plugin is started:

\begin{itemize}
\item \%D - subject of user's certificate,
\item \%P - name of credentials' proxy file.
\end{itemize}

\end{itemize}

\item \textbf{\textit{unixvo}}\textit{=vo rule} - same as
  \textbf{\textit{unixgroup}} for users belonging to Virtual
  Organization (VO) \textit{vo}.
\item \textbf{\textit{unixmap}}\textit{={[}unixname]{[}:unixgroup]
  rule} - define a local UNIX user and optionally group used to
  represent connected client. \textit{rule} is one of those allowed
  for \textbf{\textit{authorization groups}} (see
  Section~\ref{sec:Authorization}) and for
  \textbf{\textit{unixgroup}}/\textbf{\textit{unixvo}}.  In case of a
  mapping rule, username is the one provided by the rule. Otherwise
  the specified \textit{unixname:unixgroup} is taken. Both
  \textit{unixname} and \textit{unixgroup} may be either omitted or
  set to '{*}' to specify missing value.
\item \textbf{\textit{groupcfg}}\textit{=name} - is put into
  subsections representing a plugin or {[}group] section and defines
  if that section is effective. The only unaffected option is
  \textbf{\textit{groupcfg}}.  If name is empty (or no groupcfg is
  used at all), following lines apply to all users.
\end{itemize}

\subsection{Plugin Configuration}

Subsections of the \emph{gridftpd} section specify plugins which serve
the virtual FTP path (similar to the UNIX mount command). The name of
the subsection is irrelevant but it is useful to use a name related to
the plugin, e.g.\ {[}gridftpd/jobs] for the \emph{jobplugin}. Inside
the subsection, the following commands are supported:

\begin{itemize}
\item \textbf{\textit{plugin}}\textit{=library\_name} - use plugin
  \textit{library\_name} to serve virtual path.
\item \textbf{\textit{path}}\textit{=path} - virtual path to serve.
\end{itemize}

The GFS comes with 3 plugins: \textit{fileplugin.so, gaclplugin.so} and
\textit{jobplugin.so}.

\subsubsection{JobPlugin}

\textit{jobplugin.so} supports the following options:

\begin{itemize}
\item \textbf{\textit{configfile}}\textit{=path} - defines
  non-standard location of the A-REX configuration file,
\item \textbf{\textit{allownew}}\textit{=yes|no} - specifies if new
  jobs can be submitted. Default is \emph{yes}.
\item
  \textbf{\textit{unixgroup}}/\textbf{\textit{unixvo}}/\textbf{\textit{unixmap}}
  - same options as in the top-level GFS configuration. If the mapping
  succeeds, the obtained local user will be used to run the submitted
  job.
\item \textbf{\textit{remotegmdirs}}\textit{=control\_dir session\_dir
  [drain]} - specifies control and session directories under the
  control of another A-REX to which jobs can be assigned (see
  Section~\ref{sub:Multiarex}). Remote directories can be added and
  removed without restarting the GFS. However, it may be desirable to
  drain them prior to removal by adding the \emph{``drain''}
  option. In this case no new jobs will be assigned to these
  directories but their contents will still be accessible.
\item \textbf{\textit{maxjobdesc}}\textit{=size} - specifies maximal
  allowed size of job description in bytes. Default value is 5MB. If
  value is missing or set to 0 no limit is applied.
\end{itemize}

\subsubsection{FilePlugin}

\textit{fileplugin.so} supports the following options:

\begin{itemize}
\item \textbf{\textit{mount}}\textit{=path} - defines the place on
  local filesystem to which file access operations apply.
\item \textbf{\textit{dir}}\textit{=path options} - specifies access
  rules for accessing files in \textit{path} (relative to virtual and
  real path) and all the files below. \\ \textit{options} is a list of
  the following keywords:

\begin{itemize}
\item \textbf{\textit{nouser}} - do not use local file system rights, only
use those specified in this line.
\item \textbf{\textit{owner}} - check only file owner access rights.
\item \textbf{\textit{group}} - check only group access rights.
\item \textbf{\textit{other}} - check only ``others'' access rights.
\end{itemize}

The options above are exclusive. If none of the above are specified, the usual
UNIX access rights are applied.

\begin{itemize}
\item \textbf{\textit{read}} - allow reading files.
\item \textbf{\textit{delete}} - allow deleting files.
\item \textbf{\textit{append}} - allow appending files (does not allow creation).
\item \textbf{\textit{overwrite}} - allow overwriting of existing files
(does not allow creation, file attributes are not changed).
\item \textbf{\textit{dirlist}} - allow obtaining list of the files.
\item \textbf{\textit{cd}} - allow to make this directory current.
\item \textbf{\textit{create}} \textit{owner:group permissions\_or:permissions\_and}
- allow creating new files. File will be owned by \textit{owner} and
owning group will be \textit{group}. If '{*}' is used, the user/group
to which connected user is mapped will be used. The permissions will
be set to \textit{permissions\_or} \& \textit{permissions\_and} (the second
number is reserved for future usage). 
\item \textbf{\textit{mkdir}} \textit{owner:group permissions\_or:permissions\_and}
- allow creating new directories.
\end{itemize}
\end{itemize}

\subsubsection{GACLPlugin}

\textit{gaclplugin.so} supports the following options:

\begin{itemize}
\item \textbf{\textit{gacl}}\textit{=gacl} - GACL XML.
\item \textbf{\textit{mount}}\textit{=path} - local path served by
plugin.
\end{itemize}

The GACL XML may contain variables which are replaced with values taken from
the client's credentials. The following variables are supported:

\begin{list}{--}{\setlength{\labelwidth}{0.5cm}\setlength{\rightmargin}{\leftmargin}}
\item [{\emph{\$subject}}] - subject of user's certificate (DN),
\item [{\emph{\$voms}}] - subject of VOMS\cite{voms} server (DN),
\item [{\emph{\$vo}}] - name of VO (from VOMS certificate),
\item [{\emph{\$role}}] - role (from VOMS certificate),
\item [{\emph{\$capability}}] - capabilities (from VOMS certificate),
\item [{\emph{\$group}}] - name of group (from VOMS certificate) .
\end{list}
Additionally, the root directory must contain a \emph{.gacl} file with initial
ACLs. Otherwise the rule will be {}``deny all for everyone''.


\subsection{Authorization}
\label{sec:Authorization}

% this section is taken from the Config_Auth document
% ``Configuration and Authorisation of ARC (NorduGrid) Services''

ARC services which have to authorize remote client applications use
the notion of \textit{group} for authorization purposes. Each
\textit{group} is made of \textit{rules} applied sequentially. If a
client's credentials pass \textit{all rules}, the client is treated as
belonging to the specified \textit{group}.

Each group is represented by a top level section named
{[}\textbf{\textit{group}}{]} or its subsection. Each such section
represents a separate authorization group and its name is given by the
\textbf{\textit{name}} command inside that section. If there is no
\textbf{\textit{name}} command then the name of the subsection is
used.

Authorization is performed by applying set of rules. Rules obey same
format as the rest of the configuration file. Each rules command
consists of a \textit{rule word} prepended with optional
\textit{modifiers} - {[}+|-{]}{[}!{]}

\begin{itemize}
\item [+]accept credential if matches the following rule (positive match, default
action),
\item [-]reject credential if matches the following rule (negative match),
\item [!]invert matching. Match is treated as non-match. Non-match is treated
as match, either positive ({}``+'' or nothing) or negative ({}``-'').
\end{itemize}

Processing of rules in every group stops after the first positive or
negative match, or failure is reached. If a rule does not match then
processing continues.  Failures are rule-dependant and may be caused
by conditions like a missing file, unsupported rule, etc.

The following \textit{rule words} and arguments are supported:

\begin{itemize}
\item {[}\textbf{\textit{subject}}{]}\textit{=subject}
  {[}\textit{subject} {[}...{]}{]} - match user with one of specified
  subjects
\item \textbf{\textit{file}}={[}\textit{filename} {[}...{]}{]} - read
  rules from specified files (format of file is similar to Globus
  grid-mapfile with user names ignored)
\item \textbf{\textit{remote}}={[}\textit{ldap://host:port/dn}
  {[}...{]}{]} - match user listed in one of specified LDAP
  directories (uses network connection hence can take time to process)
\item \textbf{\textit{voms}}\textit{=vo group role capabilities} -
  accept user with VOMS proxy with specified \textit{vo},
  \textit{group}, \textit{role} and \textit{capabilities}. '{*}' can
  be used to accept any value
\item \textbf{\textit{vo}}={[}\textit{vo} {[}...{]}{]} - match user
  belonging to one of specified Virtual Organizations as defined in
  \emph{vo} configuration section (see below).
\item \textbf{\textit{group}}={[}\textit{groupname}
  {[}\textit{groupname} {[}...{]}{]}{]} - match user already belonging
  to one of specified groups.
\item \textbf{\textit{plugin}}\textit{=timeout plugin}
  {[}\textit{arg1} {[}\textit{arg2} {[}...{]}{]}{]} - run external
  \emph{plugin} (executable or function in shared library) with
  specified arguments. Execution of \emph{plugin} may not last longer
  than \emph{timeout} seconds.  If \emph{plugin} looks like
  \emph{function@path} then function \emph{int
    function(char{*},char{*},char{*},...)} from shared library
  \emph{path} is called (\emph{timeout} is not functional in that
  case). Rule matches if exit code is 0. In arguments following
  substitutions are applied before plugin is started:

\begin{itemize}
\item \%D - subject of user's certificate,
\item \%P - name of credentials' proxy file.
\end{itemize}

\item \textbf{\textit{lcas}}\textit{=library directory database} -
  call LCAS functions to check rule. Here \textit{library} is path to
  shared library of LCAS, either absolute or relative to
  \textit{directory}; \textit{directory} is path to LCAS installation
  directory, equivalent of LCAS\_DIR variable; \textit{database} is
  path to LCAS database, equivalent to LCAS\_DB\_FILE variable. Each
  arguments except \textit{library} is optional and may be either
  skipped or replaced with '{*}'.
\item \textbf{\textit{all}} - accept any user
\end{itemize}

Here is an example of authorization group:

\begin{verbatim}
 [group/admins]

 -subject="/O=Grid/OU=Wrong Place/CN=Bad Person"
 file="/etc/grid-security/internal-staff"
 voms="nordugrid * * admin"
\end{verbatim}

In this example the following rules are applied to determine whether
the identity presented is part of the group ``admins'':

\begin{itemize}
\item If the identity is ``/O=Grid/OU=Wrong Place/CN=Bad Person'' it
  is rejected
\item If the identity is in the mapfile
  ``/etc/grid-security/internal-staff'' the identity is accepted
\item If the identity is a VOMS proxy with a NorduGrid VO extension
  and admin capability, the identity is accepted
\end{itemize}

\subsubsection{Virtual Organizations}

VOs are defined in the \emph{vo} configuration section. The following
commands are supported:

\begin{itemize}
\item \textbf{\textit{vo}}\textit{=vo\_name} - specifies name of
  VO. Mandatory command.
\item \textbf{\textit{file}}\textit{=path} - path to file which
  contains list of users' DNs belonging to VO.
\item \textbf{\textit{source}}\textit{=URL} - specifies URL from which
  list of users may be obtained. May be multiple.
\end{itemize}


\section{Configuration Examples}

The examples presented below contain full configuration examples for
the GridFTP server. The A-REX and information system configurations
are not shown - this is described in other documents.


\subsection{Simple Example}

In the following minimal example we use a single static mapfile which
contains all possible user mappings for this site.

\begin{verbatim}
[common]
hostname="myhost.org"
lrms="fork"
gridmap="/etc/grid-security/grid-mapfile"

[gridftpd]
debug="3"
logfile="/var/log/arc/gridftpd.log"
logsize="10000000 2"
pidfile="/var/run/gridftpd.pid"
pluginpath="/usr/local/lib/arc"
encryption="no"
allowunknown="no"
maxconnections="200"

[gridftpd/jobs]
path="/jobs"
plugin="jobplugin.so"
\end{verbatim}


\subsection{Detailed Example}

Here we configure a simple PBS based cluster according to the
following use case.  John is member of the VO "smscg" where he belongs
to the group "atlas" and has been assigned the roles "production" and
"test". Since groups and roles are fully decoupled, John can request
proxies that can include one (or several) of the following different
group-role combinations (termed "Fully Qualified Names" (FQAN)):
\begin{itemize}
\item /smscg (notice it's the same as /smscg/Role=NULL)
\item /smscg/Role=production
\item /smscg/Role=test
\item /smscg/atlas
\item /smscg/atlas/Role=production
\item /smscg/atlas/Role=test
\end{itemize}

A-REX serves as front-end to a batch-system that provides a
"low\_prio\_queue" and a "high\_prio\_queue". Assignment to the
different queues is done via local user identites.  More precisely,
the local users "smscg001, smscg002, smscg003" will be assigned to the
low\_prio\_queue, whereas users "smscgP001, smscgP002, smscgP003" to
the high\_prio\_queue (the configuration of the batch-system to support
this is out of scope of this example).

Users sending jobs to A-REX should be assigned to one of the queues
depending on the credentials they present in their proxy
certificate. The assignment shall look as follows:
\begin{itemize}
\item /smscg , /smscg/Role=test , /smscg/Role=production => shall map
  to one of the smscg00[1-3] local identities (thus low\_prio\_queue)
\item /smscg/atlas , /smscg/atlas/Role=test ,
  /smscg/atlas/Role=production => shall map to one of the
  smscgP00[1-3] local identities (thus high\_prio\_queue)
\end{itemize}

The following usage pattern is considered. User John first wants to
run a monitoring job on the high\_prio\_queue. He performs a
voms-proxy-init and specifies his "/smscg/atlas/Role=test" FQAN to be
used. When he submits his monitoring-job, John will be mapped to one
of the smscgP001, smscgP002, smscgP003 accounts. John's job will thus
run on the high\_prio\_queue.

After submitting the monitoring job, John submits regular jobs with
his FQAN "/smscg". These jobs will run on the low\_prio\_queue. Later
John switches back to the FQAN "/smscg/atlas/Role=test" to fetch the
result of his monitoring job.

The discrimination to what queue John is to be mapped is done with VO
information only and not on the basis of the DN of John's
certificate. Hence the choice to what queue to be mapped is under
control of John (we silently presumed John knows the mappings at the
source).

Notes:
\begin{itemize}
\item a DN based grid-mapfile is generated on the front-end with a
  default mapping entry for John. The grid-mapfile is only used by the
  information system (GIIS) to make the grid resource look eligible
  for jobs submitted by John.
\item the DN based grid-mapfile per se does not permit John to access
  the grid resource under different local identities (e.g. once as
  smscg001 and later as smscgP001), since the first matching DN
  defines the local identity John is to be mapped to. This is not a
  flaw since NorduGrid has support for lcmaps, which allows a
  're-mapping' of a user.
\item the mapping of the FQAN to the local user identity
  (e.g. "/smscg" to local user "smscg001") shall be done with lcmaps
  (in detail the lcmaps framework + lcmaps voms plugins). Direct VOMS
  based mapping is also possible.
\end{itemize}

If user John creates a proxy certificate with the "grid-proxy-init"
command instead of "voms-proxy-init", hence the proxy certificate will
not contain any VO information and submits a job to A-REX (the
matchmaking will still work, since it's done with John's DN) he shall
not be authorized.

Example configuration:

\begin{verbatim}
[common]
pbs_bin_path="/usr/bin"
pbs_log_path="/var/spool/pbs/server_logs"
hostname="myhost.org"
lrms="pbs"


[vo]
# We will use this configuration block for a few purposes.
# 1. To generate grid-mapfile needed for information system.
#    For that purpose nordugridmap utility will have to be 
#    run periodically.
# 2. To provide coarse-grained information to authorization 
#    rules used to define authorization groups. If needed of
#    course.
id="smscg_vo"
vo="smscg_vo"

# Here we define path to file to which nordugridmap will write DNs of
# users matching rules below. Because we are going to use it as
# grid-mapfile for other purposes it is going to reside at default
# location.  
file="/etc/grid-security/grid-mapfile"

# Now we tell nordugridmap to pull information from
# VOMRS/VOMSS/or_whatever_it_is_called_now service and to ask for
# users belonging to smscg VO.
source="vomss://voms.smscg.org:8443/voms/smscg"

# Now we specify default mapping to local *NIX id. It is possible to
# completely redefine mapping in [gridftpd] block. But this one will
# be used by information system to compute and present resources
# available to user.  Let's use one of lowest priority account defined
# in use-case.
mapped_unixid="smscg001"


[group]
# In this authorization group we are going to check if user presents
# any proof that he belongs to 'smscg' VO. We can use that information 
# later to explicitly limit access to resources. If such access 
# control is not needed this group can be removed.
name="smscg_auth"

# Here we can use internal support of ARC for VOMS attributes
# voms="smscg * * *"
# If we want to limit access to resources also by other VOMS
# attributes then other voms rules similar to those defined
# below in [gridftpd] section may be used.

# Or we can ask some external executable to analyze delegated
# credentials of user. In this example executable vomatch 
# is called with first argument containing path to delegated
# proxy certificate and second - required VO name.
# plugin="10 /opt/external/bin/vomatch %P smscg"

# Or - probably preferred way in this use case - we can use 
# LCAS to analyze delegated proxy.
# First element after '=' sign is path to LCAS library whatever
# it is called in current implementation. Second is LCAS installation
# path - it will be used to set environment variable LCAS_DIR.
# And third element is path to LCAS database file - it will be passed
# to environment variable LCAS_DB_FILE.
# Function 'lcas_get_fabric_authorization' of specified LCAS library
# will be called with following 3 arguments
#  1. char* pointer to string containing DN of user
#  2. gss_cred_id_t variable pointing at delegated credentials of user
#  3. char* pointer to empty string
# Returned 0 int value is treated as positive response
lcas="/opt/glite/lib/liblcas.so /opt/glite /opt/glite/share/lcas.db"

# As coarse grained solution it is also possible to check if user 
# belongs to one of defined VOs as specified in _previously_ defined
# [vo] group. Here we refer to VO group smscg_vo defined above.
#vo="smscg_vo"


[gridftpd]
debug="2"
logfile="/var/log/arc/gridftpd.log"
logsize="10000000 2"
pidfile="/var/run/gridftpd.pid"
port="2811"
pluginpath="/usr/local/lib/arc"
encryption="no"

# By specifying 'no' here we limit users allowed to establish
# connection to this server to those specified in grid-mapfile. This
# may be not necessary if additional authorization is applied as done
# below. But this provides additional layer of protection so let it
# be.
allowunknown="no"

maxconnections="200"

# Here we start fine-grained user mapping. Let's first define few VOMS
# mappings using embedded functionality of ARC.  These lines should
# map Grid users to high-priority and low-priority *NIX users smscg001
# and smscgP001. Mind order - those with more attributes defined come
# first.  I do not know if missing attribute is passed by VOMS as
# empty string or as string containing NULL keyword. Here I assume
# empty string. If it is NULL then "" has to be replaced with NULL.
#unixmap="smscgP001 voms smscg atlas test *
#unixmap="smscgP001 voms smscg atlas production *
#unixmap="smscgP001 voms smscg atlas "" *
# These 3 lines are not needed if grid-mapfile defines default mapping
# to smscg001 user. But we can have them for consistence and if mapping
# to nobody is defined below for safety reasons.
#unixmap="smscg001 voms smscg "" test *
#unixmap="smscg001 voms smscg "" production *
#unixmap="smscg001 voms smscg "" "" *

# Instead of using multiple unixmap commands above we may define 
# 2 authorization groups using [group] blocks. Let's say their
# names are smscg_low and smscg_high. Then 'group' matching rule
# may be used. 
#unixmap="smscgP001 group smscg_high"
#unixmap="smscg001 group smscg_low"

# Or if we want to use all 6 local accounts and let mapping choose
# randomly within 2 group accounts 'simplepool' may be used. In
# example below 'unixgroup' ensures proper choice of group and
# 'simplepool' makes a choise from accounts in pool. Last argument
# specifies directory containing file named 'pool'. That file contains
# list of local user accounts. Also this directory will be used for
# writing information about current mappings.
#unixgroup="smscg_high simplepool /var/nordugrid/smscg_high"
#unixgroup="smscg_low simplepool /var/nordugrid/smscg_low"

# And mapping preferred in this use case - through LCMAPS.  First
# element after '=' sign is path to LCMAPS library whatever it is
# called in current implementation. Second is LCMAPS installation path
# - it will be used to set environment variable LCMAPS_DIR.  And third
# element is path to LCMAPS database file - it will be passed to
# environment variable LCMAPS_DB_FILE. Those 3 arguments are followed
# list of policy names.
# Function 'lcmaps_run_and_return_username' of specified LCMAPS library
# will be called with following arguments
#  1. char* pointer to string containing DN of user
#  2. gss_cred_id_t variable pointing at delegated credentials of user
#  3. char* pointer to empty string
#  4. char** pointer for chosen username.
#  5. int variable containing number of policies
#  6. char** list of policy names
# Expected 0 int value returned and argument 4 set. Value returned in
# 4th argument is used as username of local account.
unixmap="* lcmaps /opt/glite/lib/liblcmaps.so /opt/glite \
  /opt/glite/share/lcmaps.db policy1 policy2"

# Here we can specify mapping to some harmless local user account for
# safety reasons. If that account is not allowed to submit jobs to
# LRMS then this will also work as authorization effectively cutting
# off users without proper VOMS attributes.
unixmap="nobody:nobody all"


[gridftpd/jobs]
# This block defines job submission service
path="/jobs"
plugin="jobplugin.so"

# Line below specifies that this plugin/service is only available to
# users belonging to authorization group. If such behavior is not
# required then this line must be commented.
groupcfg="smscg_auth"


[queue/low_prio_queue]
name="low_prio_queue"
homogeneity="True"
scheduling_policy="FIFO"
comment="This queue is low priority"
nodecpu="adotf"
nodememory="512"
architecture="adotf"
opsys="Mandrake 8.0"
opsys="Linux-2.4.19"
benchmark="SPECINT2000 222"
benchmark="SPECFP2000 333"
cachetime="30"
timelimit="30"
sizelimit="5000"


[queue/high_prio_queue]
name="high_prio_queue"
homogeneity="True"
scheduling_policy="FIFO"
comment="This queue is high priority"
nodecpu="adotf"
nodememory="512"
architecture="adotf"
opsys="Mandrake 8.0"
opsys="Linux-2.4.19"
benchmark="SPECINT2000 222"
benchmark="SPECFP2000 333"
\end{verbatim}


\section{Running the service}

An initialization script \emph{a-rex-gridftpd} for the GFS is provided
in \emph{\$ARC\_LOCATION/etc/init.d} (or equivalent depending on
architecture).

\begin{verbatim}
Usage: a-rex-gridftpd {start|stop|status|restart|reload|condrestart}
\end{verbatim}

Upon starting and depending on the configured log level, messages will
be logged in the log file specified in the configuration file.


\section{Using Multiple A-REX services Under One GFS\label{sub:Multiarex}}

For large clusters, using a single machine for all input and output
file transfer, as well for the interaction with the LRMS and
Information System, can limit the job throughput of the
cluster. Running several A-REXs on separate hosts can help spread the
hardware and network load. A single GFS can feed jobs to several A-REXs,
hence a cluster with many A-REXs still appears as a single site to the
outside world. When a job is submitted, the GFS jobplugin assigns a
random control directory to use for the job from the main
\emph{controldir} specified in the A-REX configuration and any extra
\emph{remotegmdirs} specified in the jobplugin configuration.

Each control directory is used by a separate A-REX, therefore for every
\emph{remotegmdirs} command in the GFS configuration, there must be a
A-REX running which defines the corresponding \emph{controldir} and
\emph{sessiondir} in its configuration file. Each A-REX can run
independently on its own host, the only requirement is that the
control and session directories must be accessible on the GFS host,
and the GFS user must have write access to these directories. It is
recommended that these directories are local to the GFS host and
exported (via NFS for example) to the other hosts, rather than being
on a remote filesystem and exported to the GFS host. This means that
any glitches in the network do not cause the GFS host to hang. It is
also important that the local user accounts on each host must be
synchronised with the GFS host. A A-REX is not aware that any other A-REXs
are running, as they only see what the GFS decides should go into
their own control directory. All communication between the GFS and a
A-REX is through the A-REX's control directory. Note that in remote control
directories there is no way to specify control directories per user,
as with the \emph{control} command. Only the \emph{controldir} command
can be used and all users will use the same control directory.

One feature of this design is that multiple A-REXs can share the same
LRMS, and hence compete with each other to submit jobs. Therefore any
LRMS settings in the configuration files must be carefully matched in
order not to bias one A-REX over another. In most cases each host's
configuration file can be identical apart from the control and session
directories. Some configuration sections such as the GFS and infosys
sections will be ignored by the remote A-REX hosts as these services are
not running.

Cacheing can be set up in a variety of different ways. Each A-REX can
have its own cache, completely independent from any other, which will
lead to popular files being replicated in many caches. Or, all caches
can be shared with all A-REXs, which means no replication between
caches but heavy intra-site network traffic if the cache file systems
are hosted on different hosts. Another option is to give each A-REX
its own cache, but access to the other caches as remote caches. This
avoids replicating files and intra-site network traffic. Replication
can still be enabled by specifying ``replicate'' as the link\_path for
remote cache dirs, and the advantage of this is that files are copied
from the remote cache rather than being downloaded again from source.

When setting up an extra A-REX, it is important that no other services
(GFS, infosys) run on the host. The instances of these services
running on the ``main'' host take care of all the A-REXs. In other
words, the startup scripts for \emph{gridftpd} and \emph{grid-infosys}
should be removed from any place where they would be started
automatically (usually \emph{\$ARC\_LOCATION/etc/init.d/}). No host
certificates are required for hosts which only run a A-REX
instance. Note that in this multiple A-REX set up, there is a one to
one relationship between control and session directories, hence
multiple \emph{sessiondir} commands cannot be used in a A-REX
configuration.


\bibliography{grid}

\end{document}
